% =============================================================================
% GLOSSÁRIO HISTÓRICO E PERICIAL — TRILÍNGUE (PT | EN | ZH)
% =============================================================================
\cleardoublepage
\thispagestyle{empty}
\chapter*{Glossário Histórico e Pericial --- Historical and Forensic Glossary --- 历史与鉴定词汇表}
\addcontentsline{toc}{chapter}{Glossário · Glossary · 词汇表}
\markboth{Glossário · Glossary · 词汇表}{}

\begingroup

\noindent\textit{Vocabulário de apoio para navegar os documentos da Seca de 1915 e os laudos psiquiátricos do Hospital Colônia de Barbacena (1917--1979). --- Support vocabulary for navigating the documents of the 1915 drought and the psychiatric reports of the Barbacena Hospital-Colony (1917--1979). --- 为浏览1915年大旱文献和巴巴塞纳收容院精神科报告（1917—1979）提供的辅助词汇。}

\vspace{0.8em}
\noindent\rule{\textwidth}{0.8pt}
\vspace{0.6em}

\noindent\textbf{1. TERMOS DA SECA E DO SERTÃO (1914--1915) --- TERMS OF THE DROUGHT AND THE SERTÃO --- 旱灾与塞尔唐术语}

\vspace{0.4em}
\begin{itemize}[itemsep=0.3em]
  \item \textbf{Campo do Alagadiço --- Alagadiço Camp --- 阿拉加迪索（Alagadiço）收容营（1915）:} Campo de confinamento de retirantes criado próximo a Fortaleza durante a seca de 1915; é o referente histórico do cenário ficcional da primeira parte. --- Camp for the confinement of retirantes created near Fortaleza during the 1915 drought; it is the historical referent for the fictional setting in Part One. --- 1915年大旱期间在福塔莱萨附近建立的逃荒者收容营，是本书第一部分虚构场景的历史参照。
  \item \textbf{Campo do Patu --- Patu Camp --- 帕图（Patu）收容营（1932--1933）:} Campo situado em Senador Pompeu e associado à seca de 1932; não é o campo de 1915 narrado no percurso de Joaquim. --- Camp located in Senador Pompeu and associated with the 1932 drought; it is not the 1915 camp in Joaquim's route. --- 位于塞纳多尔·蓬佩乌，与1932年大旱相关；它并非若阿金1915年路线中的收容营。
  \item \textbf{Retirantes / Molambudos --- Retirantes / Molambudos --- 逃荒者 / 莫兰布多斯（破衣人）:} Designação popular pejorativa dos sertanejos que fugiam a pé do interior, em farrapos de pano (molambos), tentando sobreviver à seca extrema. --- Popular, pejorative designation for the sertanejos who fled the interior on foot, in rags of cloth (molambos), trying to survive the extreme drought. --- 对徒步逃离内陆、衣衫褴褛（molambos，破布）以求在极端干旱中活命的塞尔唐人的通俗贬称。
  \item \textbf{Vala Comum --- Mass Grave --- 乱葬坑:} Poço coletivo raso onde centenas de corpos de vítimas da fome e do tifo eram depositados diariamente, sem identificação, nos campos de 1915. --- Shallow collective pit where hundreds of bodies of famine and typhus victims were deposited daily, unidentified, in the concentration camps of 1915. --- 浅层集体坑穴，1915年集中营中每天都有数百具饥荒和斑疹伤寒受害者的遗体被匿名丢入其中。
\end{itemize}

\vspace{0.8em}
\noindent\rule{\textwidth}{0.5pt}
\vspace{0.6em}

\noindent\textbf{2. TERMOS MÉDICOS E PSIQUIÁTRICOS DO SÉCULO XX --- MEDICAL AND PSYCHIATRIC TERMS --- 二十世纪医学与精神科术语}

\vspace{0.4em}
\begin{itemize}[itemsep=0.3em]
  \item \textbf{Hospital Colônia de Barbacena --- Barbacena Hospital-Colony --- 巴巴塞纳收容院:} Instituição psiquiátrica fundada em Minas Gerais em 1903. Entre 1930 e 1980 tornou-se centro de internação em massa de pessoas marginalizadas sem diagnóstico psiquiátrico formal, com mais de 60.000 mortos (o chamado \textit{Holocausto Brasileiro}). --- Psychiatric institution founded in Minas Gerais in 1903. Between 1930 and 1980 it became a center of mass internment of marginalized people without formal psychiatric diagnosis, resulting in the death of more than 60,000 people (an episode known as the \textit{Brazilian Holocaust}). --- 1903年建于米纳斯吉拉斯州的精神病机构。1930至1980年间，它成为对没有正式精神科诊断的边缘人群进行大规模收容的中心，导致六万多人死亡（即所谓的\textit{巴西大屠杀}事件）。
  \item \textbf{Eletroconvulsoterapia (ECT / Eletrochoque) --- Electroconvulsive Therapy (ECT / Electroshock) --- 电休克治疗（ECT / 电休克）:} Aplicação de corrente elétrica ao cérebro para induzir convulsões. No Colônia, praticada de forma contínua e punitiva, sem anestesia ou gel condutor, causando queimaduras e fraturas. --- Application of electric current to the brain to induce seizures. At the Hospital-Colony, practiced continuously and punitively, without anesthesia or conductive gel, causing burns and fractures. --- 向大脑施加电流以诱发癫痫发作。在收容院，这种做法持续且带有惩罚性，不使用麻醉或导电凝胶，造成烧伤以及股骨和脊柱骨折。
  \item \textbf{Lobotomia Transorbital (Procedimento de Freeman) --- Transorbital Lobotomy (Freeman Procedure) --- 经眶脑叶切除术（弗里曼术式）:} Intervenção cirúrgica introduzida nos anos 1940: uma haste metálica (semelhante a um picador de gelo) inserida pelo canto interno do olho para seccionar conexões do lobo frontal. --- Surgical intervention introduced in the 1940s: a metal rod (similar to an ice pick) inserted through the inner corner of the eye to sever connections of the frontal lobe. --- 1940年代引入的外科干预，将一根金属棒（类似冰锥）从眼角内侧插入大脑，以切断额叶的连接。
  \item \textbf{Caquexia --- Cachexia --- 恶病质:} Estado de desnutrição crônica extrema e depleção muscular profunda. --- State of extreme chronic malnutrition and profound muscular wasting. --- 极端慢性营养不良和深度肌肉消耗的状态。
  \item \textbf{Linhas de Beau --- Beau's Lines --- 博氏线:} Sulcos transversais nas unhas causados pela interrupção temporária do crescimento ungueal durante episódios graves de fome ou infecção. --- Transverse grooves in the nails caused by the temporary interruption of nail growth during severe episodes of hunger or infection. --- 指甲上的横向沟纹，由严重饥饿或感染发作期间指甲生长暂时中断引起。
  \item \textbf{Icterícia Unilateral --- Unilateral Jaundice --- 单侧黄疸:} Coloração amarelada restrita a um único olho, achado pericial anômalo documentado no exame de Joaquim (Paciente 1.260). --- Yellowish coloration restricted to a single eye, an anomalous forensic finding documented in the examination of Joaquim (Patient 1,260). --- 仅限于单只眼睛的黄色着色，这是在对若阿金（患者1,260）的检查中记录到的一种异常法证发现。
\end{itemize}

\vspace{0.8em}
\noindent\rule{\textwidth}{0.5pt}
\vspace{0.6em}

\noindent\textbf{3. TERMOS PERICIAIS E ARQUIVÍSTICOS (2026) --- FORENSIC AND ARCHIVAL TERMS --- 法证与档案术语}

\vspace{0.4em}
\begin{itemize}[itemsep=0.3em]
  \item \textbf{Leucótomo --- Leucotome --- 脑白质刀:} Instrumento cirúrgico metálico pontiagudo usado em lobotomias. --- Pointed metal surgical instrument used in lobotomies. --- 用于脑叶切除术的尖头金属外科器械。
  \item \textbf{Lâmina Cribriforme --- Cribriform Plate --- 筛板:} Placa óssea perfurada no teto da cavidade nasal/ocular que dá acesso à base do cérebro. --- Perforated bone plate at the roof of the nasal/ocular cavity that provides access to the base of the brain. --- 鼻腔/眼腔顶部带孔的骨板，可通往大脑底部。
  \item \textbf{Pacto Ficcional Ergódico --- Ergodic Fictional Pact --- 遍历小说契约:} Dispositivo narrativo em que a estrutura física da obra e as decisões do leitor simulam o próprio contágio descrito no diário. --- Narrative device in which the physical structure of the work and the reader's decisions simulate the very contagion described in the diary. --- 一种叙事装置，作品的实际结构和读者的决定模拟了日记中所描述的传染本身。
\end{itemize}

\endgroup
\clearpage
